\section{Fakta, Predikat, dan Arity dalam Pemodelan}

Pemilihan predikat yang tepat beserta \textit{arity}-nya merupakan keputusan desain fundamental dalam pemodelan basis pengetahuan. Predikat yang dipilih harus mampu menangkap informasi domain secara cukup (\textit{expressive}) tanpa menjadi terlalu rumit (\textit{parsimonious}) \cite{rosen2019}.

\subsection{Predikat Arity-1: Properti}

Predikat dengan satu argumen menyatakan properti atau klasifikasi suatu entitas:

\begin{itemize}
    \item $Manusia(x)$: ``$x$ adalah manusia''
    \item $Laki(x)$: ``$x$ berjenis kelamin laki-laki''
    \item $Prima(n)$: ``$n$ adalah bilangan prima''
    \item $Login(u)$: ``pengguna $u$ telah login''
\end{itemize}

Predikat arity-1 analog dengan himpunan: $Manusia(x)$ berarti $x \in \text{Himpunan Manusia}$.

\subsection{Predikat Arity-2: Relasi Biner}

Predikat dengan dua argumen menyatakan relasi antara dua entitas:

\begin{itemize}
    \item $OrangTua(x, y)$: ``$x$ adalah orang tua dari $y$''
    \item $Suka(x, y)$: ``$x$ menyukai $y$''
    \item $LebihBesar(x, y)$: ``$x$ lebih besar dari $y$''
    \item $Teman(x, y)$: ``$x$ adalah teman $y$''
\end{itemize}

Perhatikan bahwa urutan argumen sangat penting. $OrangTua(Hasan, Ahmad)$ berbeda maknanya dari $OrangTua(Ahmad, Hasan)$.

\subsection{Predikat Arity-$n$ ($n \geq 3$): Relasi Multi-Argumen}

Untuk relasi yang melibatkan lebih dari dua entitas atau parameter, digunakan predikat dengan arity yang lebih tinggi:

\begin{itemize}
    \item $Tewas(x, penyebab, tahun)$: arity-3
    \item $Penerbangan(maskapai, asal, tujuan, waktu)$: arity-4
    \item $Nilai(mahasiswa, matakuliah, semester, skor)$: arity-4
\end{itemize}

\begin{catatan}
Dalam praktik pemodelan, arity yang terlalu tinggi (misalnya $> 5$) seringkali menandakan bahwa predikat perlu didekomposisi menjadi beberapa predikat yang lebih sederhana. Decomposisi ini mirip dengan proses normalisasi dalam desain basis data relasional.
\end{catatan}

\subsection{Fakta Ground vs Kalimat Berkuantor}

Perbedaan fundamental antara fakta dan aturan:

\begin{table}[H]
\centering
\renewcommand{\arraystretch}{1.3}
\begin{tabular}{|p{3.5cm}|p{4cm}|p{4.5cm}|}
\hline
\textbf{Aspek} & \textbf{Fakta (Ground Atom)} & \textbf{Aturan (Klausa Horn)} \\ \hline
Variabel & Tidak ada (semua konstanta) & Mengandung variabel terikat \\ \hline
Kuantor & Tidak diperlukan & Kuantor universal $\forall$ \\ \hline
Interpretasi & Kebenaran spesifik & Generalisasi / hukum \\ \hline
Contoh & $Laki(Hasan)$ & $\forall x\,(Laki(x) \rightarrow \neg Perempuan(x))$ \\ \hline
\end{tabular}
\caption{Perbandingan Fakta dan Aturan dalam KB}
\end{table}

\begin{contoh}
\textbf{Pemodelan Hobi Mahasiswa}

Fakta:
\begin{align*}
&Suka(Ali, Matematika), \quad Suka(Budi, Informatika), \\
&Suka(Cici, Fisika), \quad Suka(Ali, Informatika)
\end{align*}

Aturan:
\[ \forall x, y, z\, (Suka(x, z) \wedge Suka(y, z) \wedge (x \neq y) \rightarrow Rekan(x, y)) \]

Artinya: ``$x$ adalah rekan $y$ jika keduanya menyukai mata kuliah yang sama dan merupakan entitas yang berbeda.''

Dari fakta dan aturan di atas, kita dapat menurunkan: $Rekan(Ali, Budi)$ (karena keduanya menyukai Informatika).
\end{contoh}
